\section{Introduction}
\label{sec:intro}

% STAGE 3: rewritten 9 September 2026.  Decisions taken: RQ1--RQ3 stay; the
% second thesis is stated in its weak form (non-monotonic, not inverted U).

Large language models have reshaped automated machine learning. Systems built
on them plan experiments, write and repair code, read execution logs and
interpret their own results, which classical AutoML frameworks, searching a
fixed space of pipelines, never did. The design space of such systems is
large and largely uncharted: implementations differ in how much autonomy they
grant the model, how they inject knowledge about libraries and methods, and
how tightly they constrain the sequence of actions.

How much each of these design choices adds to accuracy remains unmeasured. Papers in the
area typically introduce one system and report that it improves on a previous
one, with the backbone model, the scaffold and the evaluation protocol all
changing at once. Comparisons across papers therefore cannot separate the architecture from the
model. Agentic systems are also rarely measured against the baseline a
competent practitioner would reach for, tuned gradient boosting; where gradient
boosting appears at all, it is usually untuned.

This paper isolates the architectural axis. We fix the backbone to a single
open 31B-parameter model and vary only the surrounding system across three
points of the design space. We call the model weak in a relative sense: it is
small against the 80B to 1.6T backbones of comparable studies
(Section~\ref{sec:related}), and two models one generation older return
nothing usable (Section~\ref{sec:modelaxis}). \textbf{FEDOT.LLM} is a rigid pipeline in which the
model configures a fixed AutoML backend through templates with fixed input and
output contracts. \textbf{AutoDS-Tools} is a multi-agent system of six
specialised agents that communicate through structured reports and have broad
freedom of action. \textbf{Terminus-2} is an open harness that gives the model
a shell, a working directory and an execution loop, and prescribes nothing
else. All three run inside one execution environment that records every
attempt, error, token and unit of cost. All are evaluated on the official
temporal splits of TabReD~\citep{rubachev2024tabred}, a benchmark of
real-world tabular tasks under distribution shift, against the eighteen
reference methods published with it. Those references are Optuna-tuned and
averaged over fifteen seeds, so they set a ceiling that reflects competent
practice. A second benchmark,
MLAgentBench~\citep{huang2024mlagentbench}, supplies a working script in every
task and so provides a second starting point.

We ask three questions.
\begin{description}
  \item[RQ1.] With the model fixed, how much does the choice of system move
        the result on real-world tabular tasks, and which expert-tuned
        baseline does the best system reach?
  \item[RQ2.] How much of that effect belongs to the architecture itself, and
        how much to what ships with it: the instruction file, the container
        image and the choice of model?
  \item[RQ3.] Does the instruction file keep its value when the task starts
        from a working script, and when the same text is handed to
        Terminus-2?
\end{description}
Section~\ref{sec:ceiling} answers RQ1 and RQ2 on TabReD,
Section~\ref{sec:mlab} answers RQ3 on MLAgentBench, and
Section~\ref{sec:mechanism} reads the execution traces for the mechanism
behind both.

The contributions are as follows.
\begin{enumerate}
\setlength{\itemsep}{0pt}
\renewcommand{\labelenumi}{(\roman{enumi})}
\item We measure how far the system around a fixed open model reaches. On the
  eight TabReD tasks, against eighteen Optuna-tuned methods, AutoDS-Tools is
  level with tuned LightGBM and $\gapXGB$ short of tuned XGBoost on a scale
  where 0 is the weakest and 1 the strongest published method per task; its
  average rank is $\rkAutoDS$ of \nMethods{}, against $\rkTerminus$ for
  Terminus-2 and $\rkFedot$ for FEDOT.LLM, which are level. No system
  approaches the strongest published method.
\item We separate the architecture from what ships with it. AutoDS-Tools
  leads Terminus-2 by $\gapAgentsShip$ on that scale; $\layerWorthNorm$ of
  that is its instruction file, the prescribed-knowledge layer of
  Section~\ref{sec:axis}, and $\gapAgentsOff$ the six-agent architecture;
  with the file removed the three architectures lie within $\spanArch$ of
  each other. The container image changes nothing measurable, and a backbone swap across
  a \axPriceSpread$\times$ price range moves accuracy by under one per cent
  at the median.
\item We split the layer into the two instructions that earlier ablations
  switch together. On TabReD, naming the library improves the task metric by
  $\gridToolMeanU\%$ on average, advice on how to train and validate by
  $\gridDiscMeanU\%$, and both together by $\gridBothMeanU\%$, below the
  $\gridSumHalvesU\%$ sum of the two.
\item We test the same layer at a different starting point and in a
  different architecture. Where every task starts from a working script
  (MLAgentBench) it leaves no submission on \mlabWiped{} of \mlabRepeated{}
  tasks and helps on one. Appended to Terminus-2 it leaves accuracy unchanged
  and loses \ktDiscLost{} of \ktDiscTrials{} trials with the training advice
  and \ktToolLost{} of \ktToolTrials{} with the library instruction; deleting
  two paragraphs about training time restores every trial. Together (iii)
  and (iv) explain why published ablations of knowledge injection disagree.
\item We check external validity on three published scientific tasks: maize
  yield, supercomputer job failure and polymer glass transition. Given the
  task text alone, Terminus-2 matched or beat the author baselines where the
  split is comparable and rediscovered the published featurisation; the text
  that names LightAutoML, on which AutoDS-Tools ran, cost Terminus-2 all
  \caseTermPrescTotal{} attempts.
\item We release every task, instruction file, trace and script behind these
  numbers.
\end{enumerate}

Across the three settings no single system wins, and the choice depends on
what the task supplies.
\begin{itemize}
\setlength{\itemsep}{0pt}
\item AutoDS-Tools is the most accurate on both benchmarks, level with tuned
  LightGBM on TabReD and ahead of Terminus-2 on \headWinsAutoDS{} of
  \headTasks{} MLAgentBench tasks, and the slowest: \minAutoDSTabred{} minutes
  per TabReD task against \minTerminusTabred{} for Terminus-2.
\item Terminus-2 is second on both benchmarks and loses more trials. On the
  scientific tasks, given the task text alone, it matched or beat the author
  baselines in one to three minutes per attempt.
\item FEDOT.LLM is level with Terminus-2 on TabReD and the only system that
  tunes its models: best of the three on one TabReD task and on F-DATA, at
  $\minFedotTabredMin$ to $\minFedotTabredMax$ minutes a task.
\end{itemize}

The design also bears on a common assumption. Terminal-Bench 2.0 introduces
Terminus-2 as a minimal scaffold, built so that a leaderboard over it reads
as a comparison between models~\citep{merrill2026terminalbench}. That reading
is sound only if the scaffold contributes little, and in the coding domain it
is already under pressure~\citep{harnessbench2026,scaffoldeffect2026}. We ask
the same question on tabular data, where expert-tuned baselines exist. On
TabReD the instruction file that AutoDS-Tools ships with moves the result by
$\layerWorthNorm$ on the normalised scale, and replacing the backbone under
the same system by one \axPriceSpread$\times$ the price moves it by $0.04$
(Sections~\ref{sec:grid} and~\ref{sec:modelaxis}). A leaderboard over one
scaffold therefore measures the scaffold's instructions as much as the model.
